%! TeX root = ../main.tex

\chapter{Introduction}

In strongly correlated materials the interplay between electrons
driven by the Coulomb interaction
can give rise to interesting phenomena, e.g.,
Kondo physics in heavy fermions~\cite{Stewart1984, Hewson1993},
the Mott metal-to-insulator transition~\cite{Belitz1994, Imada1998},
and many more~\cite{Martin2016}.
A popular method to study these materials is dynamical mean-field theory
(DMFT)~\cite{Metzner1989, Georges1996, Held2007}
which maps a system of interacting electrons on a lattice
to a system where electronic interactions are only present at an impurity
that is coupled to a non-interacting bath.
Many approaches exist to solve this model,
including quantum Monte Carlo (QMC)~\cite{Hirsch1986, Gull2011, Wallerberger2016, Wallerberger2019},
Wilson's numerical renormalization group (NRG)~\cite{Wilson1975, Bulla1998, Bulla1999, Bulla2008},
and exact diagonalization (ED)~\cite{Jaklic1994, Caffarel1994, Capone2007, Lu2014, Lu2019},
with advantages and disadvantages between all of them.
A strong argument for QMC is the fact that it can easily be extended to
impurities containing multiple orbitals which is often necessary to describe realistic materials.
However, the problem is that sampling is done on the imaginary axis which requires
analytic continuation to obtain real frequency data,
which is an ill-conditioned inversion~\cite{Jarrell1996, Beach2000, Kaufmann2023}.
In contrast, NRG --- which is designed for impurity problems ---
gives quantities on the real axis,
has excellent resolution for small frequencies
but suffers at larger ones due to logarithmic discretization~\cite{Bulla2008}.
ED methods can be applied very broadly, obtain quantities on the real axis as well
but suffer from discretization artifacts.
Continuous quantities need to be discretized to a few poles
to keep the exponentially growing many-body Hilbert space manageable.
In~\cite{Lu2014, Lu2019} it was shown how to counter this exponential increase in complexity
by finding an optimized one-particle basis in combination with a restricted active space (RAS).
However, due to discretization, the self-energy contained non-physical values.
In this work we combine their approach together with the
symmetric improved estimator for the self-energy introduced for NRG~\cite{Kugler2022}
and also used for QMC~\cite{Kaufmann2019}.
This allows us to obtain physical self-energies even for a discrete bath.
In contrast to NRG, we introduce no broadening and calculate it directly on the real axis.
With this approach we study the Mott metal-to-insulator transition
on the half-filled $z\to\infty$ Bethe lattice.

\section{Outline}

The thesis is organized as follows:
In \zcref{ch:models-and-formalism} we describe the necessary model and formalism for this work.
We start by introducing the Bethe lattice and the limit of infinite dimensions.
Afterward we create model Hamiltonians to describe interacting electrons,
introduce the Green's function formalism and the self-energy.
Finally, we explain DMFT and
arrive at a set of equations which need to be solved self-consistently.
In \zcref{ch:numerical-methods} we describe the numerical methods used to develop our code.
We explain the natural impurity orbitals~\cite{Lu2014, Lu2019} and the RAS approach,
allowing us to solve our system in polynomial instead of exponential time.
This is followed by the introduction of the Lanczos algorithm and representation of complex
functions as a sum of discrete poles.
\zcref{ch:results} focuses on the results of our work.
Finally, we summarize and discuss possible enhancements in \zcref{ch:conclusion}.
